% Source: https://artofproblemsolving.com/community/c3h2522939_geo_probability
% Stable identity: aops-topic-2522939-q01
Let $g$ be the guess and $a$ the correct answer; $(g,a)$ is uniform on the rectangle $[1,2]\times[1,3]$ of area $2$. He is correct when $|g-a|\le0.1a$, i.e. $0.9a\le g\le1.1a$, i.e.
\[\dfrac{g}{1.1}\le a\le\dfrac{g}{0.9}.\]
For $1\le g\le2$ the upper bound $\dfrac g{0.9}\le\dfrac{20}{9}<3$ stays inside the rectangle, while the lower bound $\dfrac g{1.1}$ is below $1$ when $g<1.1$. So the favourable area is
\[\int_1^{1.1}\left(\dfrac{g}{0.9}-1\right)dg+\int_{1.1}^{2}\left(\dfrac{g}{0.9}-\dfrac{g}{1.1}\right)dg=\dfrac{1}{60}+\dfrac{20}{99}\cdot\dfrac{4-1.21}{2}=\dfrac1{60}+\dfrac{31}{110}=\dfrac{197}{660}.\]

\begin{asy}
usepackage("amsmath");
size(5cm);
defaultpen(fontsize(10pt));
real sx=3, sy=1.6;
pair T(real g,real a){return (sx*(g-1),sy*(a-1));}
fill(T(1,1)--T(1,1/0.9)--T(2,2/0.9)--T(2,2/1.1)--T(1.1,1)--cycle,rgb(0.8,0.88,1));
draw(T(1,1)--T(2,1)--T(2,3)--T(1,3)--cycle,linewidth(0.8));
draw(T(1,1/0.9)--T(2,2/0.9),blue); draw(T(1.1,1)--T(2,2/1.1),blue);
label("$a=\dfrac{g}{0.9}$",T(1.45,1.45/0.9),NW,blue+fontsize(8pt));
label("$a=\dfrac{g}{1.1}$",T(1.72,1.25),blue+fontsize(8pt));
label("$1$",T(1,1),S,fontsize(8pt)); label("$2$",T(2,1),S,fontsize(8pt)); label("$1.1$",T(1.12,1),S,fontsize(7pt));
label("$1$",T(1,1),W,fontsize(8pt)); label("$3$",T(1,3),W,fontsize(8pt));
label("guess $g$",T(1.5,0.8),S,fontsize(8pt)); label(rotate(90)*"answer $a$",T(0.85,2),W,fontsize(8pt));
\end{asy}

(For the first integral: $\dfrac{1.21-1}{1.8}-0.1=\dfrac{7}{60}-\dfrac{6}{60}$.) Dividing by the total area $2$, the probability is
\[\boxed{\dfrac{197}{1320}\approx0.149}.\]
